\section{Results VII: the lockbox}\label{sec:lockbox}

All modelling choices (covariates, number of experts, penalties, refit schedules, trading rules and
their parameters) were fixed on the development sample ending 30 June 2024. We then re-estimated every
model from scratch on data up to 30 June 2024 and ran the same walk-forward through 28 October 2026,
a period of 563 to 849 trading days per series that includes the September 2024 policy rally in China
and the April 2025 tariff shock. The lockbox was evaluated once; no result in
this section fed back into any choice.

\paragraph{Density forecasts.} The ranking of model families survives intact (\cref{tab:lockbox}). The
four mixture models occupy the top four average ranks on log score (LSTAR $3.05$, SMoE $3.11$, MAR
$3.21$, deep SMoE $3.26$), and the SMoE has the best average rank on QLIKE ($2.95$) and on the FZ0 tail
loss ($3.16$). The SMoE is in the 10\% model confidence set for every one of the nineteen series under
all four losses. Pairwise, it is better than GARCH-$t$ on 14 series (significantly on 3, significantly
worse on none), and significantly better than the Markov-switching regression on 12, the Gauss-HAR on
9, the boosted model on 7 and both foundation models on 18 series. Against the other mixtures (MAR,
deep SMoE) the lockbox is too short to separate the models: each comparison produces one significant
win and one significant loss.

\paragraph{Calibration.} The SMoE's 1\% VaR passes the Kupiec test on 18 of 19 series (GARCH-$t$: 18;
Chronos-Bolt: 1). Its PIT standard deviation is $1.01$, whereas the foundation models' remain at
$1.31$--$1.35$. After adaptive conformal calibration no model fails the Kupiec test on any series, and
the SMoE and deep SMoE pass the conditional-coverage test on every series at the 5\% level.

\paragraph{Realised variance.} The 5-minute data cover only the first 246 lockbox days. On QLIKE against
realised variance the mixtures again lead (MAR $0.567$, LSTAR $0.598$, SMoE $0.606$, GARCH-$t$ $0.663$).
The Mincer--Zarnowitz $R^2$, however, is highest for the two foundation models ($0.55$ and $0.71$ against
at most $0.13$ for the others). The $R^2$ of a regression on 246 days is dominated by the realised
variance spike of late September 2024, and the zero-shot models, which condition on the raw return path,
scaled up faster after the first large moves. We report this as the one lockbox result in which a
foundation model leads; on the robust loss it does not.

\paragraph{Trading.} The strategies do not beat buy and hold in the lockbox. On the ETF, the best
configuration at 3bp is a GARCH-$t$ mean-variance rule with Sharpe ratio $0.74$ against $0.38$ for buy
and hold, but its standard error is $0.65$, its deflated Sharpe ratio $0.65$ and the probability of
backtest overfitting across the 360 configurations rises to $0.77$ (\cref{tab:overfit_lockbox}). The
SMoE-based volatility-target rules match the buy-and-hold Sharpe ratio ($0.37$ versus $0.38$) with a
smaller maximum drawdown ($-17\%$ versus $-19\%$). The lockbox thus confirms the development-period
conclusion: the predictive distributions manage risk, and the data do not support claims of timing
alpha.

\begin{table}[!htbp]\centering\scriptsize
\caption{Development period versus lockbox: average ranks across 19 series (1 = best), model-confidence-set membership (10\%) and Kupiec passes}\label{tab:lockbox}
\begin{adjustbox}{max width=\textwidth}
\begin{tabular}{llrrrrrrrr}
\toprule
Period & Model & Rank NLL & Rank CRPS & Rank QLIKE & Rank FZ5 & MCS NLL & MCS QLIKE & MCS FZ5 & Kupiec 1\% pass\\
\midrule
Development 2012--2024/06 & SMoE & 2.95 & 3.63 & 3.47 & 4.26 & 15 & 18 & 18 & 17\\
Development 2012--2024/06 & Deep-SMoE & 3.26 & 5.26 & 3.53 & 3.63 & 14 & 17 & 18 & 15\\
Development 2012--2024/06 & MAR & 3.47 & 3.42 & 3.58 & 3.95 & 15 & 18 & 18 & 16\\
Development 2012--2024/06 & LSTAR & 4.00 & 5.00 & 4.05 & 4.00 & 11 & 17 & 18 & 15\\
Development 2012--2024/06 & GARCH-t & 4.11 & 5.26 & 5.95 & 5.16 & 11 & 11 & 15 & 10\\
Development 2012--2024/06 & AR-GJR-t & 3.63 & 3.68 & 5.47 & 4.11 & 17 & 15 & 18 & 12\\
Development 2012--2024/06 & Gauss-HAR & 7.58 & 8.16 & 4.84 & 5.16 & 2 & 16 & 15 & 4\\
Development 2012--2024/06 & MS-AR & 7.89 & 9.89 & 8.00 & 7.42 & 1 & 6 & 11 & 16\\
Development 2012--2024/06 & XGB-Gauss & 8.47 & 4.63 & 7.00 & 8.26 & 1 & 10 & 9 & 2\\
Development 2012--2024/06 & Chronos-Bolt-base & 10.47 & 7.95 & 10.26 & 10.63 & 0 & 2 & 2 & 0\\
Development 2012--2024/06 & HistGauss & 10.26 & 11.47 & 10.00 & 9.68 & 1 & 6 & 7 & 1\\
Lockbox 2024/07--2026/10 & SMoE & 3.11 & 3.42 & 2.95 & 3.16 & 19 & 19 & 19 & 18\\
Lockbox 2024/07--2026/10 & Deep-SMoE & 3.26 & 3.53 & 3.89 & 3.74 & 19 & 19 & 19 & 17\\
Lockbox 2024/07--2026/10 & MAR & 3.21 & 3.11 & 3.16 & 3.58 & 19 & 19 & 19 & 18\\
Lockbox 2024/07--2026/10 & LSTAR & 3.05 & 3.32 & 3.63 & 3.68 & 19 & 19 & 19 & 18\\
Lockbox 2024/07--2026/10 & GARCH-t & 4.37 & 6.21 & 6.63 & 6.11 & 18 & 19 & 18 & 18\\
Lockbox 2024/07--2026/10 & AR-GJR-t & 5.05 & 6.16 & 6.95 & 6.11 & 18 & 19 & 18 & 17\\
Lockbox 2024/07--2026/10 & Gauss-HAR & 7.58 & 6.68 & 4.26 & 4.68 & 15 & 19 & 19 & 14\\
Lockbox 2024/07--2026/10 & MS-AR & 7.16 & 9.05 & 7.79 & 7.84 & 18 & 18 & 18 & 15\\
Lockbox 2024/07--2026/10 & XGB-Gauss & 8.32 & 4.95 & 6.37 & 6.26 & 17 & 19 & 19 & 6\\
Lockbox 2024/07--2026/10 & Chronos-Bolt-base & 10.74 & 9.58 & 10.63 & 11.05 & 4 & 8 & 9 & 1\\
Lockbox 2024/07--2026/10 & HistGauss & 10.47 & 11.32 & 10.05 & 10.05 & 12 & 18 & 16 & 8\\
\bottomrule
\end{tabular}
\end{adjustbox}
\par\smallskip\begin{minipage}{\textwidth}\footnotesize
Ranks among the twelve forecasters (Chronos-Bolt-small not shown). MCS and Kupiec columns count series out of 19.
\end{minipage}
\end{table}
\begin{table}[!htbp]\centering\small
\caption{Lockbox: backtest overfitting diagnostics (CSI 300 ETF, 360 configurations at 3bp)}\label{tab:overfit_lockbox}
\begin{adjustbox}{max width=\textwidth}
\begin{tabular}{p{6cm}rrr}
\toprule
Best at 3bp & Sharpe & DSR & PBO\\
\midrule
GARCH-t, meanvar, c3.0, g5.0, v0.8 & 0.738 & 0.648 & 0.770\\
\bottomrule
\end{tabular}
\end{adjustbox}
\end{table}

\subsection{Do foundation models remember the test period?}\label{sec:fm-release}

A zero-shot forecaster is out of sample only if its pretraining corpus ends before the forecast
date. The documented corpora of Chronos and TimesFM~2.5 do not list the series we study, but
they cannot be audited series by series, and \citet{sarkar2024lookahead} and
\citet{glasserman2024assessing} show that look-ahead through pretraining is real for language
models used in finance. A model's release date $C$ bounds its data cut-off from above, which gives a
simple test.

\begin{proposition}[Release-date test]\label{prop:release}
Let $d_t=\mathrm{NLL}^{\mathrm{FM}}_t-\mathrm{NLL}^{\mathrm{ref}}_t$ be the log-score loss of a
foundation model relative to a reference model estimated in real time, averaged across series on date
$t$, and let $d^{\mathrm{pl}}_t$ be the same difference for a placebo model that is also estimated in
real time. Suppose that, without pretraining overlap, $\E[d_t-d^{\mathrm{pl}}_t]$ is the same before
and after $C$ (parallel trends). Then the coefficient $b$ in
$d_t-d^{\mathrm{pl}}_t=a+b\,\1\{t\ge C\}+e_t$ is zero, and its Newey--West $t$-statistic is
asymptotically standard normal as the numbers of pre- and post-release dates grow. If the corpus
contains the evaluation series up to some cut-off $C'\le C$, forecasts before $C'$ are partly
in-sample and $b>0$. If $C'<C$, dates in $[C',C)$ are classified as pre-release although they are
unseen, which biases $b$ towards zero; the test remains valid but loses power.
\end{proposition}

\noindent The proof is the usual argument for a difference in means with a HAC variance; the content
is in the identifying assumption. Shifts in market conditions at $C$ that change how all models
perform relative to the SMoE are removed by the placebo, for which we use GARCH-$t$. We apply the test
with the release dates of Chronos-Bolt (26 November 2024) and TimesFM~2.5 (15 September 2025), using
the development and lockbox forecasts together.

%RESULTS-RELEASE%
